\documentclass[12pt]{article}
\usepackage[T2A]{fontenc}
\usepackage[utf8]{inputenc}
\usepackage[russian]{babel}
\usepackage{amsmath,amssymb}
\usepackage{geometry}
\geometry{margin=2.5cm}
\usepackage{indentfirst}

\title{Термодинамическая неопределённость в квантовых машинах Отто\\ с конечным временем цикла}
\author{И.~Н.~Лебедев, М.~П.~Соколова}
\date{2024}

\begin{document}
\maketitle

\section*{Аннотация}
Исследуется связь между квантовой неопределённостью и диссипацией в квантовых тепловых машинах, работающих по циклу Отто за конечное время. Выведено обобщённое соотношение неопределённости для работы и КПД, учитывающее спектр рабочего тела.

\vspace{0.5cm}
\noindent\textbf{Ключевые слова:} квантовая термодинамика; машина Отто; соотношения неопределённости; конечное время; ионные ловушки.

\vspace{0.3cm}
\noindent\textbf{Раздел журнала:} Механика. \quad \textbf{Файл:} 1 из 3.

\section{Введение}
Квантовые тепловые машины работают в условиях, где неопределённость играет существенную роль. Мы рассматриваем цикл Отто за конечное время.

\section{Основной результат}
Обобщённое соотношение неопределённости для работы $W$ и КПД $\eta$ имеет вид
\begin{equation}
\Delta W \, \Delta \eta \geq \frac{\hbar}{2} \left| \frac{\partial \eta}{\partial t} \right|.
\end{equation}

\section{Заключение}
Предложена экспериментальная схема на цепочке ионов в ловушке.

\end{document}